%% =====================================================================
%% CNRIA 2026 — 16th Conference on Research in Computer Science
%%              and its Applications, IEEE
%%
%% Title: End-to-end IoT-to-ML pipeline for hygrothermal
%%        characterization and prediction of a vernacular
%%        CEB+Typha building in a tropical-sahelian climate
%%
%% DOUBLE-BLIND SUBMISSION
%%   - Author names are anonymized (use "Anonymous Authors")
%%   - Do NOT include institution names, acknowledgments, or
%%     self-identifying references before camera-ready
%%   - Cite own prior work in third person: "as shown in [n]"
%% =====================================================================

\documentclass[conference]{IEEEtran}

%% ---- Core packages -----------------------------------------------
\usepackage[T1]{fontenc}
\usepackage[utf8]{inputenc}
\usepackage{cite}            % IEEE-style [1,2,3] citation grouping
\usepackage{amsmath}
\usepackage{graphicx}
%% Figures are looked up in several places so the same source compiles
%% locally (scripts write to ../01_python_scripts/outputs/) and on
%% Overleaf (flat upload, or a figures/ subfolder).
\graphicspath{{./}{figures/}{outputs/}{../01_python_scripts/outputs/}}
\usepackage{booktabs}        % publication-quality tables (\toprule etc.)
\usepackage{url}
\usepackage{hyperref}
\usepackage{microtype}       % improves line breaking
\usepackage{xspace}          % smart spacing after macros

%% hyperref should come LAST (except cleveref)
\hypersetup{
  colorlinks=true,
  linkcolor=black,
  citecolor=black,
  urlcolor=blue,
  pdftitle={IoT-to-ML hygrothermal pipeline for CEB+Typha building},
  pdfauthor={Anonymous Authors},
}

%% Break points let the URL wrap inside a narrow IEEE column.
\def\UrlBreaks{\do\/\do\-\do\.\do\_}
\Urlmuskip=0mu plus 1mu\relax
\newcommand{\repourl}{\url{https://anonymous.4open.science/r/2026_cnria-1AD7/}}
%% \EDB removed: it was never invoked, and expanded to a string that
%% names the host institution.
\newcommand{\ceb}{CEB+Typha\xspace}
\newcommand{\sigrat}{$\sigma_\text{int}/\sigma_\text{ext}$\xspace}

%% =====================================================================
\begin{document}

\title{End-to-end IoT-to-ML Pipeline for Hygrothermal Characterization
  and Prediction of a Vernacular \ceb Building in a Tropical-Sahelian Climate}

%% DOUBLE-BLIND: replace with author names for camera-ready
\author{\IEEEauthorblockN{Anonymous Authors}
\IEEEauthorblockA{\textit{Anonymous Institution}}}

\maketitle

% =====================================================================
\begin{abstract}
% =====================================================================
Vernacular earthen buildings represent a significant portion of the
built stock in sub-Saharan Africa, yet their hygrothermal behavior under
IoT monitoring remains largely unquantified. This paper presents an
end-to-end open-source pipeline---from raw sensor data to probabilistic
machine-learning predictions---applied to the Ecological and Digital
Building (EDB) of a West African polytechnic school. The building
envelope combines compressed-earth blocks (CEB) and \textit{typha}
bio-based insulation. A total of 2,237 measurements at a nominal
10-minute interval, collected over two windows in July--August and
November--December 2025 spanning 41.8~days at 37\,\% completeness, are
consolidated, enriched with ERA5 reanalysis data via the Open-Meteo API,
and processed through six modular pipeline layers. Three convergent indicators quantify the
envelope thermal inertia: a damping factor \sigrat of 0.54, a median
peak-temperature lag of $+1.75$\,h between exterior and interior, and a
52.5\,\% reduction in diurnal amplitude. Interior conditions satisfy the
ASHRAE~55 tropical comfort zone 67\,\% of the time in the dry season
versus only 17\,\% during the rainy season. A HistGradientBoosting
quantile model predicting temperature increments from 36 engineered
features---including CEB-aware lagged variables and Open-Meteo exogenous
inputs---achieves a mean absolute error of 0.09\,\textdegree{}C at
15\,min and 0.23\,\textdegree{}C at 1\,h, improving on persistence by 22.9\,\% (95\,\% CI [17.5, 27.8]) and
47.7\,\% ([43.2, 51.7]) respectively, against a linear autoregressive
and a seasonal naive baseline.
Permutation importance identifies the 6-hour rolling mean of exterior
temperature as the dominant exogenous predictor, ${\approx}22\times$
more important than the instantaneous exterior temperature, which
measures the envelope's thermal memory directly from data. Exterior
moisture variables, by contrast, carry no measurable weight, so this
record yields no evidence of hygrothermal coupling. Code and data that
reproduce every table and figure are available at \repourl.
\end{abstract}

\begin{IEEEkeywords}
IoT, hygrothermal monitoring, compressed earth blocks, typha insulation,
machine learning, building thermal inertia, West Africa, smart building,
time series prediction, ERA5
\end{IEEEkeywords}

% =====================================================================
\section{Introduction}
\label{sec:intro}
% =====================================================================

The buildings and construction sector was responsible for 32\,\% of
global final energy use and 34\,\% of energy-related CO$_2$ emissions
in 2023~\cite{gsr2025}. Within that total, space cooling is the
fastest-growing end-use in buildings and already accounts for about
20\,\% of the electricity they consume~\cite{iea2018cooling}, a share
driven by hot climates.
In sub-Saharan Africa, vernacular
construction techniques---particularly compressed-earth blocks
(CEB) combined with locally harvested
bio-based insulation materials---offer an energy-efficient, low-carbon
alternative to imported cement-based construction, but remain
under-instrumented and poorly characterized in the scientific
literature~\cite{morel2001}.

The emergence of low-cost IoT sensing platforms and the availability of
high-quality open reanalysis datasets (ERA5 via Open-Meteo~\cite{zippenfenig2023})
now make it feasible to monitor such buildings continuously and to
complement sparse in-situ measurements with physically consistent
exogenous variables at near-zero marginal cost. Yet, to the best of our
knowledge, no published work has combined (i)~quantitative IoT-derived
characterization of CEB thermal inertia, (ii)~probabilistic short-term
prediction from heterogeneous sensor and reanalysis inputs, and
(iii)~a fully reproducible open-source pipeline, for a vernacular earthen
building in a West African tropical-sahelian climate.

This paper fills this gap with three main contributions:

\begin{enumerate}
\item A \textbf{six-layer open-source pipeline} (consolidation,
  Open-Meteo enrichment, preprocessing, CEB-aware feature engineering,
  quantile gradient boosting, visualization) running end-to-end from
  raw CSV sensor exports to publication-ready diagnostics and
  probabilistic forecasts.

\item A \textbf{quantitative thermal signature} of the \ceb envelope,
  derived from 2,237 measurements and three convergent metrics
  (damping ratio, peak lag, amplitude reduction). Instrumented
  monitoring of compressed-earth buildings, and of typha demonstrators
  in this region, has been reported before; what is documented here is
  the dynamic response of this particular compound assembly, measured
  at building scale rather than on material specimens.

\item A \textbf{data-driven measurement of envelope thermal memory} via
  machine-learning feature importance: the 6-hour rolling mean of
  exterior temperature is the dominant exogenous predictor of interior
  temperature at $h{+}1$\,h, ${\approx}22\times$ more important than the
  instantaneous exterior temperature and ${\approx}23\times$ more than
  its 1-hour mean. The envelope responds to an integral of outdoor
  conditions over several hours rather than to instantaneous forcing---a
  time constant recovered without physical supervision, converging with
  the measured damping and lag.
\end{enumerate}

The case study is the Ecological and Digital Building (EDB)
of a polytechnic school in Senegal,
built with CEB walls and \textit{typha} insulation within the
PNEEB-Typha programme---the Senegalese national programme for
greenhouse-gas emission reduction through energy efficiency in the
building sector, whose Typha component develops bio-based thermal
insulation from the invasive reed \textit{Typha
australis}~\cite{niang2018}. The site experiences
two climatically contrasting regimes: a humid rainy season
(June--October, RH 60--80\,\%) and a dry season with
harmattan continental winds (November--May, RH 30--55\,\%, high diurnal
amplitude).

The remainder of this paper is organized as follows.
Section~\ref{sec:related} reviews related work on earthen-building
monitoring and data-driven thermal prediction.
Section~\ref{sec:method} describes the instrumentation, data, and
pipeline architecture. Section~\ref{sec:results} presents the
hygrothermal characterization and prediction results.
Section~\ref{sec:discussion} discusses the physical interpretation,
limitations, and perspectives. Section~\ref{sec:conclusion} concludes.

% =====================================================================
\section{Related Work}
\label{sec:related}
% =====================================================================


\textbf{IoT monitoring of earthen and vernacular buildings.}
Compressed-earth construction has been studied from a material science
perspective for its thermal mass and hygroscopic
properties~\cite{morel2001,liuzzi2018}, and wall-level hygrothermal
transfer models exist for rammed earth~\cite{soudani2017}. However,
continuous IoT monitoring of complete vernacular buildings at system
level---capturing the coupled T/RH/air-quality dynamics under real
occupancy and real climate---remains sparse, particularly in West Africa.
Existing studies in the region focus either on static thermal
characterization of individual CEB specimens~\cite{fabbri2014} or on
qualitative surveys of indoor comfort without sensor data. The present
work differs by providing a longitudinal IoT time series spanning two
distinct seasonal regimes, enabling dynamic characterization of envelope
performance from the building-scale response rather than material
specimens.

\textbf{Data-driven thermal prediction for buildings.}
Machine learning has been extensively applied to building energy
forecasting in commercial, well-instrumented settings, from classical
time-series methods~\cite{deb2017} to deep architectures~\cite{li2021}.
Depth is not always required: on windowed forecasting problems a
properly engineered gradient-boosting model matches or outperforms
state-of-the-art deep networks across nine datasets~\cite{elsayed2021},
while remaining far cheaper to fit---which is what recommends it in the
low-data regime studied here. Yet the vast majority
of published models assume rich multi-zone sensor networks (10--100+
sensors), continuous multi-year records, and climate zones with mild
seasonal transitions. The low-data single-sensor regime---a single IoT
node, two discontinuous measurement windows totaling 5 months, a
single-building sample---is essentially unstudied for sub-Saharan earthen
buildings. This regime is precisely the one that will prevail in most
African research and practitioner contexts for the next decade.

\textbf{ERA5 reanalysis as exogenous input.}
Reanalysis products such as ERA5 and ERA5-Land~\cite{muñozsabater2021}
provide physically consistent gridded reconstructions of past atmospheric
conditions at 0.1\textdegree{} spatial resolution and hourly temporal
resolution, freely accessible via the Open-Meteo
API~\cite{zippenfenig2023}. Their use as exogenous model inputs has been
demonstrated for solar forecasting and agricultural applications, but their
systematic integration into building-scale thermal prediction
pipelines---particularly for gap-filling missing local sensors and
supplying dew-point, shortwave radiation, and wind speed features---has
not, to our knowledge, been reported for vernacular African buildings. Our
pipeline demonstrates that ERA5 enrichment is feasible (MAE vs.\ on-site
sensor: ${\approx}1.0$--$1.5$\,\textdegree{}C, $r>0.9$) and that it
enables a feature importance analysis otherwise impossible with a
single-variable on-site measurement. It also lets that analysis return a
negative verdict on the exogenous moisture variables
(Section~\ref{ssec:importance}), which is itself informative.

% =====================================================================
\section{Materials and Methods}
\label{sec:method}
% =====================================================================

\subsection{Study Site and Instrumentation}

The Ecological and Digital Building (EDB) is a single-story
demonstration building located on the campus of a polytechnic school in
Senegal. Exact coordinates are withheld for double-blind review and will
be released with the camera-ready version. Its envelope consists of
load-bearing compressed-earth block
(CEB) walls with \textit{typha} (\textit{Typha australis}) insulation
panels, a construction system developed and tested in the national
PNEEB-Typha programme~\cite{niang2018}. The site lies in the Sahelian
transition zone, characterized by two sharply contrasting regimes: a rainy
season (June--October, RH 60--80\,\%,
$T$ 25--32\,\textdegree{}C) and a dry season with harmattan continental
winds (November--May, RH 30--55\,\%, large diurnal amplitudes).

Instrumentation comprises a single connected sensor node (identifier
``node\_01'' in the released data) recording at a nominal 10-minute
interval. Three of its channels are used here: interior air temperature
$T_\text{int}$ (\textdegree{}C), interior relative humidity
$\text{RH}_\text{int}$ (\%), and near-envelope exterior temperature
$T_\text{ext}$ (\textdegree{}C). The node also logs an NH$_3$
air-quality index and a derived heat index; neither enters this study,
and both are dropped at load. All timestamps are in UTC.

\subsection{Data Consolidation and Quality}

Two CSV exports were provided: \texttt{sensor\_export\_1.csv} (2,237 rows) and
\texttt{sensor\_export\_2.csv} (1,548 rows, confirmed as a strict subset).
After deduplication on the \texttt{Date} field, the consolidated dataset
contains 2,237 observations spanning 16~July to 4~December 2025. Two
data quality issues were addressed. First, 391 $T_\text{ext}$ values are
absent in July, attributed to delayed sensor deployment; these are filled
by ERA5 reanalysis (Section~\ref{ssec:era5}). Second, the derived heat
index correlates with $T_\text{int}$ and $\text{RH}_\text{int}$ at
$r>0.93$; it is dropped at load rather than filtered later, so that it
cannot reach any feature set. The dataset contains a 99-day gap
(19~August--26~November 2025), separating the rainy-season window
(July--August) from the dry-season window (November--December). This
discontinuity motivates segment-aware preprocessing and precludes
long-horizon cross-season modeling.

\subsection{Data Completeness}
\label{ssec:completeness}

The record is substantially sparser than its calendar extent suggests,
and Table~\ref{tab:completeness} states this explicitly. The two seasonal
windows span 41.8~days in total, but 2,237 samples at a nominal
10-minute interval amount to only 15.5~days of data, i.e.\ an overall
completeness of 37.2\,\%. The two windows are moreover of very unequal
quality. The dry-season window is a single uninterrupted 8.0-day block
at 90.4\,\% completeness, whereas the rainy-season window is not a
continuous window at all: it consists of 12 disjoint fragments totalling
10.2~days of coverage inside a 33.8-day extent, at 24.4\,\%
completeness. Exterior temperature is sparser still in that window
(16.4\,\%), the 391 values missing in July being concentrated there.

Two consequences follow, and we state them rather than qualify them.
First, after resampling to the 15-minute modelling grid, only 41.2\,\%
of the $T_\text{int}$ and $\text{RH}_\text{int}$ slots and 33.5\,\% of
the $T_\text{ext}$ slots are backed by an actual sample; the remainder
is interpolated or, for $T_\text{ext}$ in July, substituted from
reanalysis. Second, the seasonal statistics of
Section~\ref{sec:results}---including the comfort fractions---rest on
these unequal supports, and sensor outages in the rainy season cannot be
assumed missing at random. The dry-season figures characterise eight
days at the very start of a November--May dry season, not the season as
a whole, and the rainy-season figures rest on a fragmented quarter of
their nominal window. All seasonal comparisons below should be read
under this constraint.

%% Generated by 01_python_scripts/completeness.py — regenerate rather
%% than edit by hand; the script also writes outputs/completeness.csv.
\begin{table}[!t]
\centering
\caption{Data completeness per variable and per seasonal window,
  against the nominal 10-minute sampling interval.}
\label{tab:completeness}
\begin{tabular}{@{}lrrrr@{}}
\toprule
& \multicolumn{2}{c}{Rainy (16 Jul--19 Aug)}
& \multicolumn{2}{c}{Dry (26 Nov--4 Dec)} \\
\cmidrule(lr){2-3}\cmidrule(lr){4-5}
Variable & $n$ & \% & $n$ & \% \\
\midrule
$T_\text{int}$              & 1{,}190 & 24.4 & 1{,}047 & 90.4 \\
$\text{RH}_\text{int}$      & 1{,}190 & 24.4 & 1{,}047 & 90.4 \\
$T_\text{ext}$              &     799 & 16.4 & 1{,}047 & 90.4 \\
\midrule
Window extent (days)   & \multicolumn{2}{c}{33.8} & \multicolumn{2}{c}{8.0} \\
Contiguous fragments   & \multicolumn{2}{c}{12}   & \multicolumn{2}{c}{1}   \\
Coverage (days)        & \multicolumn{2}{c}{10.2} & \multicolumn{2}{c}{8.0} \\
\bottomrule
\end{tabular}
\end{table}

\subsection{ERA5 Enrichment via Open-Meteo}
\label{ssec:era5}

Missing and supplementary meteorological variables are retrieved from the
ERA5-Land reanalysis~\cite{muñozsabater2021} via the Open-Meteo Historical
Archive API~\cite{zippenfenig2023}, using the official
\texttt{openmeteo-requests} Python library (FlatBuffers protocol,
persistent SQLite HTTP cache, automatic retry). Seven hourly variables are
fetched at the EDB coordinates: $T_\text{2m}$, $\text{RH}_\text{2m}$,
dew-point temperature, shortwave radiation, cloud cover, wind speed at
10\,m, and precipitation. The ERA5 hourly series are linearly interpolated
to the 15-minute pipeline time step and aligned to the sensor index.
Validation against the available on-site $T_\text{ext}$ measurements
yields MAE$\,{\approx}\,1.2$\,\textdegree{}C and Pearson $r>0.92$,
consistent with published ERA5-Land accuracy in West
Africa~\cite{muñozsabater2021}.

\subsection{Pipeline Architecture}

The processing pipeline is organized in six independent, testable layers
(Fig.~\ref{fig:pipeline}), implemented as four Python scripts and a batch
orchestrator. \textbf{L1 Ingestion} performs multi-file consolidation,
schema validation, deduplication, and UTC normalization.
\textbf{L2 Enrichment} fetches Open-Meteo data, validates against
on-site measurements, and aligns to the sensor index.
\textbf{L3 Preprocessing} resamples to 15\,min, segments seasonally
(rainy vs.\ dry season), winsorises outliers, and
interpolates contextual gaps. \textbf{L4 Feature engineering} produces
36 predictors (Section~\ref{ssec:features}).
\textbf{L5 Modeling} trains quantile gradient boosting and runs
walk-forward validation (Section~\ref{ssec:model}).
\textbf{L6 Visualization} generates multi-panel dashboards, calendar
heatmaps, permutation importance plots, and an automated markdown
diagnostic report. All scripts use fixed random seeds
(\texttt{random\_state=42}), relative paths, and the \texttt{matplotlib}
\texttt{Agg} backend for cross-platform headless rendering. The
repository at \repourl contains the two raw sensor exports, the
scripts, a pinned dependency list and a table mapping each command to
the table or figure it produces; running the five commands it lists
regenerates every quantitative result in this paper on Windows, macOS
or Linux.

\begin{figure*}[!t]
  \centering
  %% regen_pipeline.py → outputs/fig_pipeline.pdf (vector; a PNG
  %% preview is written too, and pdflatex picks the PDF ahead of it)
  \includegraphics[width=\textwidth]{fig_pipeline}
  \caption{End-to-end pipeline architecture. Six independently testable
    layers process the data from raw IoT sensor CSV to publication-ready
    diagnostics. The dashed branch indicates the ERA5/Open-Meteo
    Historical Archive API (L2) supplying both gap-filled exterior
    temperature and six exogenous meteorological variables to the feature
    engineering layer.}
  \label{fig:pipeline}
\end{figure*}

\subsection{CEB-Aware Feature Engineering}
\label{ssec:features}

Thirty-six predictors are constructed, exploiting physical knowledge of
the \ceb system. They include the current interior temperature and
humidity $y(t)$: withholding them costs 11--27\,\% of accuracy at short
horizons, since persistence has access to that value. Lagged exterior
temperature values at
$\Delta t \in \{15\,\text{min},\,1\,\text{h},\,2\,\text{h},\,6\,\text{h}\}$
are deliberately calibrated around the empirically measured thermal lag
(${\approx}1.75$\,h, Section~\ref{ssec:thermal}). Temporal cycles are
encoded as $(\sin,\cos)$ pairs for hour-of-day and day-of-year,
preserving continuity at midnight and year boundaries---critical for
tree-based models. Rolling statistics (mean, standard deviation) of
$T_\text{ext}$ over 1\,h and 6\,h windows capture the short memory of
the envelope. Autoregressive lags of $T_\text{int}$ and
$\text{RH}_\text{int}$ at
$\Delta t \in \{15\,\text{min},\,1\,\text{h},\,24\,\text{h}\}$ anchor
predictions to recent observed state. The six Open-Meteo exogenous
variables (dew-point, shortwave radiation and its 1\,h rolling mean,
humidity, cloud cover, wind) complete the feature set.

\subsection{Model and Validation Strategy}
\label{ssec:model}

The forecasting model is a \texttt{HistGradientBoostingRegressor}
(scikit-learn~\cite{pedregosa2011}), trained with quantile loss at
$\alpha \in \{0.05,\,0.50,\,0.95\}$ to produce both point forecasts
(median) and 90\,\% nominal prediction intervals. Hyperparameters are:
400 estimators, learning rate 0.05, maximum depth~6. Three independent
models per target variable per horizon are trained in parallel.

Validation follows a strict walk-forward strategy:
\texttt{TimeSeriesSplit} with $k=4$ expanding folds ensures no future
information leaks into training. Four prediction horizons are evaluated
independently: 15\,min, 1\,h, 6\,h, and 24\,h. The naive seasonal
baseline is the same-time previous-day value:
$\hat{y}(t) = y(t - 24\,\text{h})$, a strong baseline for smooth
periodic signals. Feature importance is assessed by permutation on a
held-out test set ($n=800$ samples), with 5 repeats per feature.

Two properties of this setup were verified rather than assumed. First,
the absence of temporal leakage: every feature at time $t$ was
recomputed on the series truncated at $t$ and compared with its value on
the full series. All 36 features are identical at every probe, so no
feature reads past $t$; the lagged and rolling terms are trailing by
construction, and the targets $y(t{+}h)$ lie strictly outside their
support. Second, the uncertainty attached to each reported gain: with
only four folds, a mean is a weak summary, so the per-sample absolute
errors of the model and of each baseline are pooled across the fold test
sets and resampled jointly (paired bootstrap, 2,000 replicates) to give
a 95\,\% interval on the difference in MAE.

% =====================================================================
\section{Results}
\label{sec:results}
% =====================================================================

\subsection{Hygrothermal Overview}

Table~\ref{tab:stats} summarizes the seasonal statistics. The interior
mean temperature is 29.6\,\textdegree{}C during the rainy season and
26.7\,\textdegree{}C in the dry season, reflecting both the lower diurnal
cycle amplitude in the dry season and the effective thermal buffering of
the CEB mass. Interior RH drops sharply between seasons: 67.0\,\% on
average during the rainy season against 45.9\,\% during the dry season.
The calendar heatmap (Fig.~\ref{fig:heatmap}) visualizes the abrupt regime
transition around late November, where cooler nights and significantly
drier air replace the sustained warm-humid conditions of the rainy season.

%% Generated by 01_python_scripts/completeness.py
%% → outputs/seasonal_summary.tex
\begin{table}[!t]
\centering
\caption{Summary statistics by seasonal regime}
\label{tab:stats}
\begin{tabular}{@{}lrrrr@{}}
\toprule
Regime     & $n$     & $T_\text{int}$ (\textdegree{}C)
           & $\text{RH}_\text{int}$ (\%) & Comfort (\%) \\
\midrule
Global     & 2{,}237 & $28.2\pm2.6$ & $57.1\pm14.8$ & 40.5 \\
Rainy      & 1{,}190 & $29.6\pm2.2$ & $67.0\pm7.9$  & 17.3 \\
Dry season & 1{,}047 & $26.7\pm2.1$ & $45.9\pm12.7$ & 66.9 \\
\bottomrule
\end{tabular}
\end{table}

\begin{figure*}[!t]
  \centering
  %% regen_heatmap.py → outputs/heatmap_calendaire_broken.png
  %% (same script also writes a continuous-axis variant,
  %%  heatmap_calendaire.png, legible but mostly empty)
  \includegraphics[width=0.68\textwidth]{heatmap_calendaire_broken}
  \caption{Calendar heatmap (hour$\,\times\,$day) of interior
    temperature, exterior temperature, and interior RH\@. The axis is
    broken at the 99-day interruption, marked on the facing edges: the
    left block spans 16~July--19~August and the right block
    26~November--4~December, drawn to a common daily scale with widths
    proportional to their durations. Within each block the calendar is
    complete, so shorter outages---such as 7--18~August---also appear at
    true width, and grey marks days without data rather than a value on
    the colour scale. Each row uses one colour scale across both blocks,
    so the two seasons remain directly comparable. Coverage is 30 of
    142~days for the interior variables and 22 for exterior temperature
    (Section~\ref{ssec:completeness}). The rainy-to-dry transition is
    visible as a simultaneous shift to cooler nights and lower humidity.}
  \label{fig:heatmap}
\end{figure*}

\subsection{Thermal Signature of the CEB+Typha Envelope}
\label{ssec:thermal}

Three convergent metrics characterize the envelope's dynamic thermal
behavior (Fig.~\ref{fig:thermal}). Damping is reported throughout as the
ratio of standard deviations \sigrat $=0.54$, computed on the 1,846
concurrent $(T_\text{int},T_\text{ext})$ samples: the interior standard
deviation, 2.50\,\textdegree{}C, is 54\,\% of the exterior value of
4.62\,\textdegree{}C. The median peak lag of $+1.75$\,h corresponds to the
delay between daily $T_\text{ext}$ and $T_\text{int}$ maxima, computed
on the 15-minute grid after a centred 1-hour rolling mean, over $n=9$
days with at least 18\,h of data (interquartile range $+1.00$ to
$+3.00$\,h, full range $+0.50$ to $+3.25$\,h). Peak picking on an
hourly grid, as in an earlier version of this analysis, quantises every
daily estimate to a whole hour and inflates the median to $+2.0$\,h; a
cross-correlation estimate over the whole record independently returns
$+1.75$\,h.

The third metric is the reduction in diurnal amplitude: over the
20~days with usable daily extrema, mean interior amplitude is
5.19\,\textdegree{}C against 10.92\,\textdegree{}C outdoors, a ratio of
0.48 and hence a 52.5\,\% reduction (Fig.~\ref{fig:thermal}, panel~c).
The least-squares fit through the same points has a slope of 0.30 with a
positive intercept of 1.88\,\textdegree{}C ($r=0.77$); it is shown for
completeness but is not a damping measure, since the intercept absorbs
part of the attenuation. Only \sigrat is used as the damping indicator
in this paper.

The all-season linear regression of interior on exterior temperature
yields $T_\text{int} = 0.42\,T_\text{ext} + 16.6$, providing a compact
empirical transfer function exploitable in simplified energy models.

The observed lag of ${\approx}2$\,h is short compared with the values
usually associated with massive earth walls, but the comparison is only
meaningful at equal wall thickness: for rammed earth, reported time lags
scale from 2.8\,h at 110\,mm to 6.2\,h at 220\,mm and 9.3\,h at
330\,mm~\cite{soebarto2009,fabbri2014}. Two mechanisms are consistent
with a lag of this order. The typha layer, being of low density and low
conductivity, adds thermal resistance without storing energy and so
shortens the transit time of the thermal wave; a moderate envelope
thickness would produce a similar effect on its own. Separating the two
requires envelope construction data (thickness, density, position of the
typha layer) that the present campaign did not record, so the
attribution to the typha layer is offered as a hypothesis rather than an
established result.

\begin{figure*}[!t]
  \centering
  %% Generated by 01_python_scripts/regen_damping.py
  %% → outputs/amortissement_CEB.png
  \includegraphics[width=0.78\textwidth]{amortissement_CEB}
  \caption{Thermal signature of the \ceb envelope: (a) 7-day sample
    showing visible delay and attenuation; (b) daily peak lag, median
    $+1.75$\,h, computed on the 15-minute grid; (c) interior vs.\
    exterior diurnal amplitude, annotated with the damping ratio
    \sigrat $=0.54$. The least-squares line in (c) is shown for
    completeness and is a different statistic from the damping ratio.}
  \label{fig:thermal}
\end{figure*}


\subsection{Hygrothermal Comfort}

Using the ASHRAE~55 tropical comfort zone adapted to naturally ventilated
buildings (24--30\,\textdegree{}C, 30--70\,\% RH~\cite{ashrae55}),
comfort time fractions differ sharply by season: 66.9\,\% in the dry
window versus only 17.3\,\% in the rainy window. These two figures
rest on very unequal supports---8~contiguous days at 90.4\,\%
completeness against 12~fragments at 24.4\,\%
(Section~\ref{ssec:completeness})---so they are reported as observed
fractions over the available samples, not as season-wide estimates. In
the dry window, mean $T_\text{int}=26.7$\,\textdegree{}C and
RH$\,{\approx}\,46$\,\% place the measurements inside the comfort
polygon. The rainy-season shortfall
is driven primarily by humidity: mean $\text{RH}_\text{int}=67$\,\% sits
just below the 70\,\% comfort ceiling, but frequent exceedances during
precipitation events push measurements outside the zone. This highlights
hygric management---natural cross-ventilation, vapor-permeable
finishes---as the priority lever for rainy-season comfort, rather than
thermal inertia alone.

\subsection{Predictive Model Performance}

\textbf{Baselines and scoring protocol.} A forecast at a 15-minute
horizon must be measured against persistence, $\hat{y}(t{+}h) = y(t)$,
which is the strong competitor on a smooth signal. We therefore report
three baselines: persistence; a linear autoregressive model refitted on
each training fold from the interior lags at 15\,min, 1\,h and 24\,h;
and the seasonal naive predictor $\hat{y}(t{+}h) = y(t{+}h{-}24\,\text{h})$.
All three are scored on exactly the same walk-forward test folds as the
model, which is what makes the comparison meaningful.

\textbf{Increment formulation.} The model predicts the increment
$y(t{+}h) - y(t)$ and the level is reconstructed by adding $y(t)$; all
metrics are computed on the reconstructed level. This matters more than
any hyper-parameter. Regressing the level directly requires boosted
trees to approximate the identity $y(t{+}h) \approx y(t)$ through
piecewise-constant splits, and the residual of that approximation is
coarser than the 0.12\,\textdegree{}C that separates two consecutive
15-minute samples. Under the level formulation the model reaches only
0.475\,\textdegree{}C MAE at 15\,min, five times worse than persistence;
under the increment formulation it reaches 0.092\,\textdegree{}C, and
persistence becomes the null prediction of zero that the model must beat.

\textbf{Results.} Table~\ref{tab:results} reports the outcome. The model
has the lowest MAE at seven of the eight horizon/target pairs, but the
bootstrap intervals show that not all of those margins are resolved.
Interior temperature carries the result: MAE 0.09\,\textdegree{}C at
15\,min, a gain of $+22.9\,\%$ over persistence
(95\,\% CI $[+17.5,\,+27.8]$); 0.23\,\textdegree{}C at 1\,h,
$+47.7\,\%$ $[+43.2,\,+51.7]$; and 0.94\,\textdegree{}C at 6\,h,
$+56.4\,\%$ $[+52.3,\,+60.3]$. All three intervals exclude zero by a
wide margin.

Humidity is a different matter. The gains over persistence are
$+4.9\,\%$ at both 15\,min $[+0.7,\,+8.9]$ and 1\,h $[-0.4,\,+9.9]$:
the first barely clears zero, the second does not clear it at all.
Persistence is already close to the sensor's noise floor on
$\text{RH}_\text{int}$, and on this record the model does not
demonstrably improve on it at short horizons. At 6\,h the humidity gain
becomes resolved again, $+12.1\,\%$ $[+5.8,\,+18.4]$.

At 24\,h neither target is resolved in either direction: temperature
loses $4.6\,\%$ to persistence with a CI of $[-10.6,\,+0.9]$ and
humidity gains $2.8\,\%$ with $[-1.8,\,+7.3]$. Persistence and the
seasonal baseline coincide at that horizon, since $y(t)$ and
$y(t{+}h{-}24\,\text{h})$ are then the same value. The honest reading is
that one day ahead the model and simple repetition are
indistinguishable on this record; the 99-day interruption is the
plausible cause, as no fold contains a continuous stretch long enough to
learn a cross-season daily cycle.

\textbf{Interval calibration.} Empirical coverage of the nominally
90\,\% intervals is 77\,\% and 71\,\% at 15\,min, and 51\,\% and 58\,\%
at 1\,h, degrading to 13--30\,\% at the longest horizons. The increment
formulation improves coverage substantially---the level formulation gave
26--42\,\% throughout---but the intervals remain over-confident. Post-hoc
conformal calibration is therefore still required before any operational
use (Section~\ref{sec:discussion}).

Per-fold results vary considerably: on temperature at 15\,min the MAE
across the four expanding folds ranges from 0.18 to 1.03\,\textdegree{}C
under the level formulation. Folds that straddle the regime change are
systematically the worst, which is a direct consequence of the
discontinuity documented in Section~\ref{ssec:completeness} rather than
of the estimator.

%% Generated from 01_python_scripts/outputs/resultats_v2_openmeteo.csv
%% (increment formulation). Per-fold values: outputs/resultats_par_fold.csv
\begin{table}[!t]
\centering
\caption{Walk-forward validation over 4 expanding folds. All columns are
  MAE on the reconstructed level, scored on identical test folds.
  \textbf{Bold} = the model beats all three baselines.}
\label{tab:results}
\setlength{\tabcolsep}{4.2pt}
\begin{tabular}{@{}llrrrrr@{}}
\toprule
& & \multicolumn{1}{c}{Model} & \multicolumn{3}{c}{Baselines}
  & \multicolumn{1}{c}{Cov.} \\
\cmidrule(lr){3-3}\cmidrule(lr){4-6}\cmidrule(lr){7-7}
Horizon & Target & MAE & Persist. & AR & Seas. & 90\,\% \\
\midrule
15\,min & T (\textdegree{}C) & \textbf{0.09} &  0.12 &  0.16 &  1.02 & 77\% \\
15\,min & RH (\%)            & \textbf{0.65} &  0.67 &  1.29 & 11.06 & 71\% \\
1\,h    & T (\textdegree{}C) & \textbf{0.23} &  0.44 &  0.36 &  1.03 & 51\% \\
1\,h    & RH (\%)            & \textbf{2.25} &  2.36 &  3.30 & 11.34 & 58\% \\
6\,h    & T (\textdegree{}C) & \textbf{0.94} &  2.16 &  1.98 &  1.01 & 36\% \\
6\,h    & RH (\%)            & \textbf{9.00} & 10.25 & 17.93 & 12.32 & 30\% \\
24\,h   & T (\textdegree{}C) &         1.05  &  0.99 &  1.08 &  0.99 & 13\% \\
24\,h   & RH (\%)            & \textbf{13.76}& 14.20 & 15.36 & 14.20 & 23\% \\
\bottomrule
\end{tabular}
\end{table}

\subsection{Feature Importance: Envelope Thermal Memory}
\label{ssec:importance}

Permutation importance on the $T_\text{int}$ model at $h{+}1$\,h
(Fig.~\ref{fig:importance}, $n=800$ held-out samples, 5 repeats)
separates the predictors sharply. Two variables dominate: the interior
autoregressive lag at 15\,min ($0.394 \pm 0.035$ decrease in $R^2$) and
the \emph{6-hour rolling mean} of exterior temperature
($0.364 \pm 0.035$). No other feature reaches a tenth of these values;
the third, the hour-of-day sine encoding, contributes $0.076$.

The informative comparison is within the exterior-temperature family
itself. Instantaneous $T_\text{ext}$ ranks fifth at $0.017$, and the
1-hour rolling mean seventh at $0.016$: the 6-hour mean is
${\approx}22\times$ more important than the instantaneous value and
${\approx}23\times$ more than the 1-hour mean. Point lags of
$T_\text{ext}$ at 1\,h and 2\,h carry no positive importance at all.
The model therefore does not respond to outdoor temperature at any
single past instant, but to its integral over roughly six hours.

This is a data-driven measurement of the envelope's thermal memory, and
it converges with the two physical indicators of
Section~\ref{ssec:thermal}: a damping ratio of 0.54 and a peak lag of
$+1.75$\,h both describe a wall that low-pass filters the outdoor
signal. The importance ranking puts a time constant on that filter, at a
scale of several hours, obtained without any physical supervision. For
practitioners the operational consequence is direct: a reduced-order
model of such an envelope needs a multi-hour aggregate of outdoor
temperature as an input, not the current reading.

We also report a negative result, since it bears on a hypothesis that
motivated the exogenous enrichment in the first place. Hygroscopic
coupling in earthen walls would predict that outdoor moisture content
helps forecast interior temperature, clay conductivity varying with
moisture state~\cite{liuzzi2018}. It does not, here: exterior dew-point
ranks 20th of 36 features with an importance of $-5\times10^{-6}$,
statistically indistinguishable from zero, and none of the six
Open-Meteo moisture or radiation variables carries measurable weight.
The enrichment remains valuable for gap-filling $T_\text{ext}$
(Section~\ref{ssec:era5}), but on this dataset it yields no evidence of
hygrothermal coupling. Whether that reflects the physics or the limited
support of the record---42~days at 37\,\% completeness, with the rainy
season, when moisture effects would be strongest, covered at only
24\,\%---cannot be settled here.

\begin{figure*}[!t]
  \centering
  %% Generated by pipeline_EDB_v2.py → outputs/v2_importance_features.png
  \includegraphics[width=0.72\textwidth]{v2_importance_features}
  \caption{Permutation importance (top-15 of 36 features) for
    $T_\text{int}$ at $h{+}1$\,h, as the decrease in $R^2$ over 5
    repeats on 800 held-out samples; bars show the standard deviation.
    Colours: red = autoregressive; purple = exterior temperature;
    green = Open-Meteo exogenous; orange = temporal cycles. The 6-hour
    rolling mean of exterior temperature (rank~2) is
    ${\approx}22\times$ more important than the instantaneous value
    (rank~5), which is the signature of the envelope's multi-hour
    thermal memory. Open-Meteo moisture and radiation variables sit at
    the bottom of the ranking; exterior dew-point, not shown, ranks
    20th of 36.}
  \label{fig:importance}
\end{figure*}

% =====================================================================
\section{Discussion}
\label{sec:discussion}
% =====================================================================

\textbf{Physical interpretation of the thermal signature.}
The measured damping ratio of 0.54 falls within the range reported in the
literature for compressed earth blocks (typically 0.5--0.7 depending on
wall thickness and density~\cite{liuzzi2018,fabbri2014}), confirming that
the CEB mass provides effective amplitude attenuation. The peak lag of
$+1.75$\,h is at the low end of the values reported for earth walls,
which range from under 3\,h for thin sections to above 9\,h for
sections of 330\,mm and beyond~\cite{soebarto2009,soudani2017}. One
reading is that the typha layer, as a low-density, low-conductivity,
low-mass material, reduces the transit time of the thermal wave without
storing energy, so that the assembly behaves as a hybrid of resistive
(typha) and capacitive (CEB) layers---a regime that classical
lumped-capacity earth models do not capture. A thin envelope would,
however, yield a comparable lag without any resistive layer, and the
present data cannot separate the two. Establishing the hybrid
interpretation requires the envelope build-up and a wall-surface or
interstitial measurement, neither of which was available here. What the
campaign does contribute is a building-scale record of this specific
compound assembly, which we have not found reported elsewhere.

\textbf{Long-horizon performance ceiling.}
The model beats all three baselines up to 6\,h but loses to persistence
on temperature at 24\,h (Table~\ref{tab:results}). This does not indicate
a failure of the model class but a structural limitation of the data: the
99-day gap between the rainy and dry seasons means no training fold
contains a continuous stretch long enough to learn a cross-season daily
cycle, so simple repetition is hard to beat one day ahead. A continuous
campaign of $\geq$6 months across at least one full seasonal transition
is the single most impactful improvement to lift this ceiling. Pre-trained
time-series foundation models (Chronos-T5~\cite{ansari2024}, Lag-Llama,
TimesFM) deployed in zero-shot are an alternative direction once denser
data becomes available.

\textbf{Calibration of prediction intervals.}
Empirical coverage of the 90\,\% nominal intervals reaches 71--77\,\% at
15\,min but falls to 13--30\,\% at 24\,h. Part of the deficit was an
artefact of predicting the level rather than the increment, which alone
cost some 30 percentage points of coverage at short horizons; what
remains is the well-known over-confidence of quantile gradient boosting
on small datasets. Conformal prediction~\cite{angelopoulos2021} provides
a finite-sample-guaranteed correction without retraining: the MAPIE
library implements it as a drop-in calibration layer. Adding this is the
priority extension and a prerequisite to any operational deployment.


\textbf{Generalizability and limitations.}
The pipeline runs unchanged on Windows, macOS, and Linux, requires only
freely available data (sensor CSV + Open-Meteo), and is directly
applicable to any earthen or vernacular building equipped with a single
connected sensor and located within ERA5 coverage---essentially the
entire African continent. The main limitations are (i)~the single-sensor
setup, which precludes characterization of intra-building thermal
gradients; (ii)~the temporal discontinuity discussed above; and
(iii)~the absence of heat-flux and air-velocity measurements, which
prevents separation of conductive, convective, and radiative
contributions to the energy balance. A second instrumentation campaign
addressing these three points constitutes the direct continuation of
this work.

% =====================================================================
\section{Conclusion}
\label{sec:conclusion}
% =====================================================================

This paper presented an end-to-end open-source pipeline for the
hygrothermal characterization and probabilistic prediction of a vernacular
\ceb building in a tropical-sahelian climate. From 2,237 IoT measurements
enriched with ERA5 reanalysis, three convergent indicators quantify the
envelope thermal inertia---damping ratio \sigrat $=0.54$, median peak
lag $+1.75$\,h, diurnal amplitude reduction 52.5\,\%---an empirical
signature of this specific compound assembly at building scale. A
HistGradientBoosting quantile model predicting increments achieves
0.09\,\textdegree{}C MAE at 15\,min and 0.23\,\textdegree{}C at 1\,h,
improving on persistence by 22.9\,\% and 47.7\,\% respectively, with
bootstrap intervals excluding zero. Permutation importance identifies the
6-hour rolling mean of exterior temperature as the dominant exogenous
predictor---${\approx}22\times$ the instantaneous value---which measures
the envelope's thermal memory directly from data and converges with the
damping and lag obtained by physical analysis. Exterior moisture
variables carry no measurable weight on this record, so the hygroscopic
coupling hypothesis is not supported by the present data.

The pipeline is released at \repourl, with the raw exports and pinned
dependencies, and is intentionally lightweight to serve as a foundation
for African research institutions instrumenting their own vernacular
buildings. The next steps---conformal calibration, multi-zone
instrumentation, BIM/ontology integration---define a clear roadmap toward
a deployable digital twin of the EDB, and more broadly toward a
standardized data-driven characterization of African earthen construction.


% =====================================================================
\bibliographystyle{IEEEtran}
\bibliography{refs}

\end{document}
